% Generated draft: requires mathematical and source review.
\documentclass[11pt]{article}
\usepackage[margin=25mm]{geometry}
\usepackage{fontspec}
\setmainfont{Latin Modern Roman}
\usepackage{amsmath,amssymb,mathtools,unicode-math}
\setmathfont{Latin Modern Math}
\usepackage{enumitem,needspace,xurl,hyperref}
\setlength{\parindent}{0pt}
\setlength{\parskip}{5pt}
\setlength{\emergencystretch}{3em}
\newcommand{\N}{\mathbb{N}}
\newcommand{\Z}{\mathbb{Z}}
\newcommand{\Q}{\mathbb{Q}}
\newcommand{\R}{\mathbb{R}}
\newcommand{\C}{\mathbb{C}}
\newcommand{\F}{\mathbb{F}}
\newcommand{\A}{\mathbb{A}}
\newcommand{\PP}{\mathbb{P}}
\newcommand{\E}{\mathbb{E}}
\newcommand{\eps}{\varepsilon}
\newcommand{\dd}{\,\mathrm{d}}
\begin{document}
\section*{UnsolvedMath ID: 30001541}
\textbf{Source problem code:} catalogue entry
\subsection{Definitions and mathematical statement}
A number field is a finite extension \(K/\Q\). An embedding \(\sigma:K\hookrightarrow\C\) is nonreal if \(\sigma(K)\not\subset\R\). A real place means an embedding into \(\R\); a field may have both real and nonreal embeddings, but a chosen real embedding cannot satisfy the displayed nonreal condition.

For a subgroup \(\Gamma\leq\operatorname{SL}_2(\C)\), its trace field is the embedded field
\[
 k_{\mathrm{tr}}(\Gamma)=\Q\bigl(\{\operatorname{tr}(\gamma):\gamma\in\Gamma\}\bigr)\subseteq\C.
\]
If the holonomy of an orientable hyperbolic \(3\)-manifold \(M\) is represented by a specified lift \(\widetilde\rho:\pi_1(M)\to\operatorname{SL}_2(\C)\), this convention gives \(k_{\mathrm{tr}}(M,\widetilde\rho)=k_{\mathrm{tr}}(\widetilde\rho(\pi_1(M)))\). The source does not specify this convention or distinguish the trace field from the invariant trace field, whose standard lifted-group expression is \(\Q(\{\operatorname{tr}(\gamma^2):\gamma\in\Gamma\})\).

\textbf{Question as supplied.} Is every pair \((K,\sigma)\) with \(\sigma(K)\not\subset\R\) the trace field of some hyperbolic \(3\)-manifold?

Under the interpretation that \(K\) is a number field, \(\sigma:K\hookrightarrow\C\) is the specified embedding, and the ordinary trace-field convention above is intended, the realization sought is
\[
\begin{gathered}
\text{for every number field }K\text{ and embedding }\sigma:K\hookrightarrow\C
\text{ with }\sigma(K)\not\subset\R,\\
\text{there exist a hyperbolic }3\text{-manifold }M
\text{ and a lifted holonomy }\widetilde\rho
\text{ such that }k_{\mathrm{tr}}(M,\widetilde\rho)=\sigma(K).
\end{gathered}
\]
This is a conditional interpretation, not a resolved formulation: the supplied question does not specify completeness, finite volume, orientability, connectedness, or treatment of boundary, nor does it specify the precise holonomy and trace-field conventions. Equality here is equality of subfields of \(\C\), recording the chosen embedding rather than only the abstract isomorphism class of \(K\).

\textbf{Unresolved source conflict.} The supplied literature assessment instead describes a number field with a real place and a chosen real embedding, and calls the pair admissible without defining admissibility. A chosen real embedding contradicts the explicit condition \(\sigma(K)\not\subset\R\). The two descriptions cannot be adopted simultaneously; the intended domain and manifold conventions require human source review.
\subsection{Short English statement}
The displayed question asks whether every number field with a specified nonreal complex embedding can occur as an embedded trace field of a hyperbolic \(3\)-manifold. The accompanying assessment instead specifies a real embedding, so the intended formulation remains unresolved. The supplied open label is dated curation, not an independently verified status.
\subsection{Sources}
\begin{itemize}
\item Supplied local registry record: permanent UnsolvedMath ID \(30001541\), title \emph{Trace Fields of Hyperbolic 3-Manifolds}, collection \emph{Oberwolfach Reports Open Problems}, exact source problem code \texttt{catalogue entry}. The supplied proposed year is 2010; the proposer field is empty. The background and literature assessment attribute the question to Neumann.
\item Supplied source citation: \emph{Low-Dimensional Topology and Number Theory} (2011), Oberwolfach Reports, pp.\ 2101--2163, DOI \texttt{10.4171/owr/2010/35}. Original OWR report: \url{https://doi.org/10.4171/owr/2010/35}. The background describes these reports as documenting mathematical workshops and their open-problem discussions.
\item Supplied literature source, labeled ``trace-field computations'': \url{https://arxiv.org/abs/1402.5582}. No authors, title, or precise supporting result are supplied for this link.
\item Literature checked at: \texttt{2026-08-21}. This is supplied dated curation, not independently verified. Its assessment labels the question open and states that later work computes fields in many families but gives no general realization theorem. That assessment uses the conflicting real-embedding formulation described above. The supplied background explicitly requests independent verification of time-sensitive status and cited claims.
\item The registry labels the statement \texttt{corrected\_verified}, explaining that the embedding condition was restored to standard mathematical notation without changing the question. This is a supplied verification label, not independent mathematical verification of this draft or resolution of the conflicting assessment.
\end{itemize}
\end{document}
